\documentclass{article}
\usepackage{amsmath,amssymb,amsthm}
\usepackage{algorithm}
\usepackage{algpseudocode}
\usepackage{hyperref}
\usepackage{enumitem}
\newtheorem{theorem}{Theorem}
\newtheorem{lemma}{Lemma}
\newtheorem{assumption}{Assumption}
\newcommand{\E}{\mathbb{E}}
\newcommand{\R}{\mathbb{R}}
\newcommand{\norm}[1]{\left\lVert #1\right\rVert}
\begin{document}

\title{SAGA for Non‑Convex Smooth Optimization \\ Detailed Algorithmic Description and Convergence Proof}
\author{}
\date{}
\maketitle

\tableofcontents
\newpage

%%%%%%%%%%%%%%%%%%%%%%%%%%%%%%%%%%%%%%%%%%%%%%%%%%%%%%%%%%%%
\section*{Algorithm Name}
\addcontentsline{toc}{section}{Algorithm Name}
\begin{center}
\textbf{SAGA (Stochastic Average Gradient Augmented) for Non‑Convex Smooth Objectives}
\end{center}
%%%%%%%%%%%%%%%%%%%%%%%%%%%%%%%%%%%%%%%%%%%%%%%%%%%%%%%%%%%%

\section*{Mathematical Formulation}
\addcontentsline{toc}{section}{Mathematical Formulation}
We consider the finite‑sum problem
\[
\min_{x\in\R^{d}} f(x)\quad\text{with}\quad 
f(x)=\frac{1}{n}\sum_{i=1}^{n}f_i(x),
\]
where each component $f_i:\R^{d}\to\R$ is differentiable (possibly non‑convex) and $L$–smooth:
\[
\forall x,y\in\R^{d},\qquad 
\|\nabla f_i(x)-\nabla f_i(y)\|\le L\|x-y\|.
\tag{A1}
\]

For each index $i$ we keep a \emph{reference point} $\phi_i\in\R^{d}$ and its stored gradient
\[
g_i:=\nabla f_i(\phi_i).
\]
Initially $\phi_i$ may be set to the starting iterate $x_0$ (or any other point).

At iteration $k$ we sample an index $i_k$ uniformly from $\{1,\dots,n\}$.
Define the \emph{SAGA gradient estimator}
\[
v_k \;:=\; \nabla f_{i_k}(x_k)-\nabla f_{i_k}(\phi_{i_k})\;+\;\frac{1}{n}\sum_{j=1}^{n} g_j .
\tag{1}
\]
The update rule is
\[
x_{k+1}=x_k-\alpha\,v_k,
\tag{2}
\]
where $\alpha>0$ is a constant stepsize.

After the update we replace the stored point for the sampled index:
\[
\phi_{i_k}\;\leftarrow\;x_k,\qquad 
g_{i_k}\;\leftarrow\;\nabla f_{i_k}(x_k).
\tag{3}
\]
All other $\phi_j$, $g_j$ ($j\neq i_k$) remain unchanged.

The algorithm repeats (2)–(3) for $k=0,1,\dots,T-1$.
%%%%%%%%%%%%%%%%%%%%%%%%%%%%%%%%%%%%%%%%%%%%%%%%%%%%%%%%%%%%

\section*{Pseudocode}
\addcontentsline{toc}{section}{Pseudocode}
\begin{algorithm}[H]
\caption{SAGA for Non‑Convex Smooth Objectives}
\begin{algorithmic}[1]
\State \textbf{Input:} stepsize $\alpha>0$, number of iterations $T$,
initial point $x_0\in\R^{d}$.
\State Initialize $\phi_i\gets x_0$ and $g_i\gets\nabla f_i(x_0)$ for all $i=1,\dots,n$.
\State Compute $\bar g\gets \frac1n\sum_{j=1}^{n}g_j$.
\For{$k=0$ \textbf{to} $T-1$}
    \State Sample $i_k\sim\text{Unif}\{1,\dots,n\}$.
    \State $v_k\gets \nabla f_{i_k}(x_k)-g_{i_k}+\bar g$.\hfill\Comment{Equation (1)}
    \State $x_{k+1}\gets x_k-\alpha v_k$.\hfill\Comment{Equation (2)}
    \State $g_{i_k}\gets \nabla f_{i_k}(x_k)$,\;
          $\phi_{i_k}\gets x_k$.\hfill\Comment{Equation (3)}
    \State Update $\bar g\gets \bar g+\frac{1}{n}\bigl(g_{i_k}^{\text{new}}-g_{i_k}^{\text{old}}\bigr)$.
\EndFor
\State \textbf{Output:} $\{x_k\}_{k=0}^{T}$ (or any randomly chosen iterate).
\end{algorithmic}
\end{algorithm}
%%%%%%%%%%%%%%%%%%%%%%%%%%%%%%%%%%%%%%%%%%%%%%%%%%%%%%%%%%%%

\section*{Dry Run Example}
\addcontentsline{toc}{section}{Dry Run Example}
Consider the simple scalar problem
\[
f(w)= (w-3)^2,
\]
which can be written as a single‑component sum $f_1(w)= (w-3)^2$, i.e. $n=1$.
We set $w_0=0$, stepsize $\alpha=0.1$, and run three iterations.

\begin{center}
\begin{tabular}{c|c|c|c}
\textbf{Iter.} & $w_k$ & $v_k=\nabla f_1(w_k)-\nabla f_1(\phi_1)+\bar g$ & $f(w_{k+1})$\\ \hline
0 & $0$ &
$\displaystyle \underbrace{2(w_0-3)}_{-6}
-\underbrace{2(\phi_1-3)}_{-6}
+\underbrace{2(\phi_1-3)}_{-6}
=-6$ & $f(0.6)= (0.6-3)^2=5.76$\\[4pt]
1 & $0.6$ &
$\displaystyle 2(0.6-3)-2(\phi_1-3)+2(\phi_1-3) = -4.8$ & $f(0.6+0.48)=f(1.08)=3.6864$\\[4pt]
2 & $1.08$ &
$\displaystyle 2(1.08-3)-2(\phi_1-3)+2(\phi_1-3) = -3.84$ & $f(1.08+0.384)=f(1.464)=2.3755$\\
\end{tabular}
\end{center}
Explanation of the first iteration (Iter.~0):
\begin{itemize}
\item $\nabla f_1(w_0)=2(w_0-3)=-6$.
\item $\phi_1$ is initialised to $w_0$, thus $\nabla f_1(\phi_1)=-6$.
\item $\bar g=\frac1n g_1 =-6$.
\item Hence $v_0=-6$ and $w_1=w_0-\alpha v_0 =0-0.1(-6)=0.6$.
\item After the update we set $\phi_1\gets w_0=0$ and $g_1\gets\nabla f_1(w_0)=-6$ (no change).
\end{itemize}
The same computation repeats because $n=1$, showing the deterministic gradient descent behaviour with stepsize $0.1$.

%%%%%%%%%%%%%%%%%%%%%%%%%%%%%%%%%%%%%%%%%%%%%%%%%%%%%%%%%%%%

\section*{Assumptions}
\addcontentsline{toc}{section}{Assumptions}
\begin{assumption}[Smoothness]\label{as:smooth}
Each component $f_i$ is $L$‑smooth:
\[
\|\nabla f_i(x)-\nabla f_i(y)\|\le L\|x-y\|\quad\forall x,y\in\R^{d}.
\]
\end{assumption}
\begin{assumption}[Bounded variance of stored gradients]\label{as:var}
Define the gradient variance at iteration $k$ as
\[
\Delta_k:=\frac{1}{n}\sum_{i=1}^{n}\|\nabla f_i(\phi_i^k)-\nabla f_i(x_k)\|^{2}.
\]
We assume $\Delta_k$ is finite for all $k$ (it will be controlled by the algorithm).
\end{assumption}
\begin{assumption}[Stepsize]\label{as:stepsize}
The constant stepsize satisfies
\[
0<\alpha\le \frac{1}{3L}.
\]
\end{assumption}
No convexity is required; the analysis holds for general non‑convex $f$.

%%%%%%%%%%%%%%%%%%%%%%%%%%%%%%%%%%%%%%%%%%%%%%%%%%%%%%%%%%%%

\section*{Main Convergence Theorem}
\addcontentsline{toc}{section}{Main Convergence Theorem}
\begin{theorem}\label{thm:main}
Let $\{x_k\}_{k\ge0}$ be generated by SAGA with stepsize $\alpha$ satisfying Assumption~\ref{as:stepsize}.
Define the Lyapunov function
\[
\Phi_k:= f(x_k)+\frac{c}{n}\sum_{i=1}^{n}\|\nabla f_i(\phi_i^k)-\nabla f_i(x_k)\|^{2},
\]
where $c:=\frac{\alpha L}{1-\alpha L}>0$.
Then for every $T\ge1$,
\[
\frac{1}{T}\sum_{k=0}^{T-1}\E\!\big[\|\nabla f(x_k)\|^{2}\big]
\;\le\;
\frac{2\big(f(x_0)-f^{\star}\big)}{\alpha T}
\;+\;
\frac{4L\alpha}{1-\alpha L}\, \E\!\big[\Delta_0\big],
\tag{4}
\]
where $f^{\star}=\inf_{x}f(x)$.
Consequently, choosing $\alpha =\Theta\!\big(\frac{1}{\sqrt{T}}\big)$ yields
\[
\min_{0\le k<T}\E\!\big[\|\nabla f(x_k)\|^{2}\big]=\mathcal{O}\!\Big(\frac{1}{\sqrt{T}}\Big).
\]
\end{theorem}
%%%%%%%%%%%%%%%%%%%%%%%%%%%%%%%%%%%%%%%%%%%%%%%%%%%%%%%%%%%%

\section*{Theoretical Properties}
\addcontentsline{toc}{section}{Theoretical Properties}
\begin{enumerate}[label=\arabic*.]
\item \textbf{Stability.} The Lyapunov function $\Phi_k$ is non‑increasing in expectation:
\[
\E[\Phi_{k+1}]\le \E[\Phi_k]-\frac{\alpha}{2}\E[\|\nabla f(x_k)\|^{2}],
\]
provided $\alpha\le 1/(3L)$. This guarantees that the iterates do not diverge even though $f$ may be non‑convex.

\item \textbf{Hyperparameter Sensitivity.}
The bound \eqref{4} contains the factor $\frac{4L\alpha}{1-\alpha L}$ multiplying the initial gradient table error $\Delta_0$.
Hence a too large $\alpha$ inflates this term, while a too small $\alpha$ slows the $1/(\alpha T)$ term.
The sweet spot $\alpha\asymp 1/\sqrt{T}$ balances both contributions.

\item \textbf{Comparison to Baselines.}
\begin{itemize}
\item \emph{SGD} with stepsize $\eta_t\propto 1/\sqrt{t}$ attains $\mathcal{O}(1/\sqrt{T})$ on the same class of problems but with a larger variance term.
\item \emph{SVRG} (epoch‑based variance reduction) also reaches $\mathcal{O}(1/T)$ for non‑convex objectives but requires periodic full‑gradient passes.
\item SAGA achieves a comparable rate \emph{online} (per‑iteration memory of $O(nd)$) and the bound \eqref{4} does not depend on the epoch length.
\end{itemize}
\end{enumerate}
%%%%%%%%%%%%%%%%%%%%%%%%%%%%%%%%%%%%%%%%%%%%%%%%%%%%%%%%%%%%

\section*{Discussion}
\addcontentsline{toc}{section}{Discussion}
The theorem shows that, despite the lack of convexity, the variance‑reduced estimator $v_k$ drives the expected squared gradient norm to zero at the same $\mathcal{O}(1/\sqrt{T})$ rate as classic SGD, but with a smaller constant thanks to the control of the table error $\Delta_k$. The Lyapunov analysis is crucial: it couples the optimisation error $f(x_k)-f^{\star}$ with the stored‑gradient discrepancy, allowing us to absorb the stochastic error introduced by the random index selection.

%%%%%%%%%%%%%%%%%%%%%%%%%%%%%%%%%%%%%%%%%%%%%%%%%%%%%%%%%%%%

\section*{Preliminaries (Definitions and Lemmas)}
\addcontentsline{toc}{section}{Preliminaries (Definitions and Lemmas)}
\begin{lemma}[Smoothness inequality]\label{lem:smooth}
If $f_i$ is $L$‑smooth then for any $x,y\in\R^{d}$,
\[
f_i(y)\le f_i(x)+\langle\nabla f_i(x),y-x\rangle+\frac{L}{2}\|y-x\|^{2}.
\]
\end{lemma}
\begin{proof}
Standard consequence of $L$‑smoothness; see e.g. Nesterov (2004). \hfill $\square$
\end{proof}

\begin{lemma}[Unbiasedness of the SAGA estimator]\label{lem:unbiased}
Conditioned on the sigma‑field $\mathcal{F}_k$ generated by $\{x_0,\dots,x_k,\phi_i^k\}_{i=1}^{n}$,
\[
\E\!\big[\,v_k\mid\mathcal{F}_k\big]=\nabla f(x_k).
\]
\end{lemma}
\begin{proof}
Because $i_k$ is uniform,
\[
\E\!\big[\nabla f_{i_k}(x_k)\mid\mathcal{F}_k\big]=\frac1n\sum_{i=1}^{n}\nabla f_i(x_k)=\nabla f(x_k).
\]
Similarly,
\[
\E\!\big[\nabla f_{i_k}(\phi_{i_k}^k)\mid\mathcal{F}_k\big]=\frac1n\sum_{i=1}^{n}\nabla f_i(\phi_i^k).
\]
Hence
\[
\E[v_k\mid\mathcal{F}_k]=\nabla f(x_k)-\frac1n\sum_{i=1}^{n}\nabla f_i(\phi_i^k)
+\frac1n\sum_{i=1}^{n}\nabla f_i(\phi_i^k)=\nabla f(x_k).
\] \hfill $\square$
\end{proof}
%%%%%%%%%%%%%%%%%%%%%%%%%%%%%%%%%%%%%%%%%%%%%%%%%%%%%%%%%%%%

\section*{Main Proof (Step‑by‑Step Derivation)}
\addcontentsline{toc}{section}{Main Proof (Step‑by‑Step Derivation)}
We prove Theorem~\ref{thm:main}. All expectations are taken w.r.t. the random index $i_k$ and the history up to iteration $k$.
\begin{enumerate}
\item[Step 1.] \textbf{Smoothness descent for $f$.} Using Lemma~\ref{lem:smooth} with $y=x_{k+1}=x_k-\alpha v_k$,
\[
f(x_{k+1})\le f(x_k)-\alpha\langle\nabla f(x_k),v_k\rangle
+\frac{L\alpha^{2}}{2}\|v_k\|^{2}.
\tag{5}
\]

\item[Step 2.] \textbf{Take conditional expectation.}
Applying $\E[\cdot|\mathcal{F}_k]$ to \eqref{5} and using Lemma~\ref{lem:unbiased},
\[
\E\!\big[f(x_{k+1})\mid\mathcal{F}_k\big]\le
f(x_k)-\alpha\|\nabla f(x_k)\|^{2}
+\frac{L\alpha^{2}}{2}\E\!\big[\|v_k\|^{2}\mid\mathcal{F}_k\big].
\tag{6}
\]

\item[Step 3.] \textbf{Bound the second moment of $v_k$.}
Write $v_k$ as
\[
v_k=\big(\nabla f_{i_k}(x_k)-\nabla f_{i_k}(\phi_{i_k})\big)
+\underbrace{\frac1n\sum_{j=1}^{n}\nabla f_j(\phi_j)}_{\bar g_k}.
\]
Hence
\[
\|v_k\|^{2}
=\big\|\underbrace{\nabla f_{i_k}(x_k)-\nabla f_{i_k}(\phi_{i_k})}_{a_k}
+\bar g_k\big\|^{2}
=\|a_k\|^{2}+2\langle a_k,\bar g_k\rangle+\|\bar g_k\|^{2}.
\tag{7}
\]

\item[Step 4.] \textbf{Take expectation of \eqref{7}.}
Because $i_k$ is independent of $\bar g_k$,
\[
\E\!\big[a_k\mid\mathcal{F}_k\big]=\frac1n\sum_{i=1}^{n}
\big(\nabla f_i(x_k)-\nabla f_i(\phi_i)\big).
\]
Thus
\[
\E\!\big[\langle a_k,\bar g_k\rangle\mid\mathcal{F}_k\big]
=\Big\langle\frac1n\sum_{i=1}^{n}
\big(\nabla f_i(x_k)-\nabla f_i(\phi_i)\big),\;\bar g_k\Big\rangle .
\tag{8}
\]

Moreover, using $L$‑smoothness,
\[
\| \nabla f_i(x_k)-\nabla f_i(\phi_i)\|^{2}
\le L^{2}\|x_k-\phi_i\|^{2}.
\tag{9}
\]

\item[Step 5.] \textbf{Introduce the Lyapunov term.}
Define
\[
\Delta_k:=\frac{1}{n}\sum_{i=1}^{n}\|\nabla f_i(\phi_i)-\nabla f_i(x_k)\|^{2}.
\tag{10}
\]
From \eqref{9} we have $\Delta_k\le L^{2}\frac1n\sum_i\|x_k-\phi_i\|^{2}$.

\item[Step 6.] \textbf{Bound $\E[\|v_k\|^{2}\mid\mathcal{F}_k]$ in terms of $\|\nabla f(x_k)\|^{2}$ and $\Delta_k$.}
A direct expansion using \eqref{7}–\eqref{8} yields
\[
\E\!\big[\|v_k\|^{2}\mid\mathcal{F}_k\big]
\le 2\|\nabla f(x_k)\|^{2}+2\Delta_k.
\tag{11}
\]
\emph{Derivation of (11):}
\begin{align*}
\E[\|v_k\|^{2}\mid\mathcal{F}_k]
&= \E[\|a_k+\bar g_k\|^{2}\mid\mathcal{F}_k] \\
&= \E[\|a_k\|^{2}\mid\mathcal{F}_k]+2\langle\E[a_k\mid\mathcal{F}_k],\bar g_k\rangle+\|\bar g_k\|^{2}\\
&\le \E[\|a_k\|^{2}\mid\mathcal{F}_k]+2\Big\|\E[a_k\mid\mathcal{F}_k]\Big\|\|\bar g_k\|
+\|\bar g_k\|^{2} \\
&\le \E[\|a_k\|^{2}\mid\mathcal{F}_k]+2\Big\|\frac1n\sum_i(\nabla f_i(x_k)-\nabla f_i(\phi_i))\Big\|
\|\bar g_k\|+\|\bar g_k\|^{2}.
\end{align*}
Now use the inequality $\|u+v\|^{2}\le 2\|u\|^{2}+2\|v\|^{2}$ with
$u=\nabla f(x_k)$ and $v=\frac1n\sum_i(\nabla f_i(\phi_i)-\nabla f_i(x_k))$ to obtain
\[
\|\bar g_k\|^{2}\le 2\|\nabla f(x_k)\|^{2}+2\Delta_k,
\qquad
\Big\|\frac1n\sum_i(\nabla f_i(x_k)-\nabla f_i(\phi_i))\Big\|^{2}\le \Delta_k.
\]
Plugging these bounds into the previous line and simplifying gives (11). \hfill $\square$

\item[Step 7.] \textbf{Plug (11) into the descent inequality \eqref{6}.}
\[
\E\!\big[f(x_{k+1})\mid\mathcal{F}_k\big]
\le f(x_k)-\alpha\|\nabla f(x_k)\|^{2}
+\frac{L\alpha^{2}}{2}\big(2\|\nabla f(x_k)\|^{2}+2\Delta_k\big).
\]
Simplify:
\[
\E\!\big[f(x_{k+1})\mid\mathcal{F}_k\big]
\le f(x_k)-\alpha\big(1-L\alpha\big)\|\nabla f(x_k)\|^{2}
+L\alpha^{2}\Delta_k.
\tag{12}
\]

\item[Step 8.] \textbf{Evolution of the table error $\Delta_k$.}
When index $i_k$ is sampled, $\phi_{i_k}$ is replaced by $x_k$.
Hence for each $i$,
\[
\|\nabla f_i(\phi_i^{k+1})-\nabla f_i(x_{k+1})\|^{2}
=
\begin{cases}
\|\nabla f_i(x_k)-\nabla f_i(x_{k+1})\|^{2}, & i=i_k,\\[4pt]
\|\nabla f_i(\phi_i^{k})-\nabla f_i(x_{k+1})\|^{2}, & i\neq i_k .
\end{cases}
\]
Using $L$‑smoothness,
\[
\|\nabla f_i(x_k)-\nabla f_i(x_{k+1})\|^{2}
\le L^{2}\|x_k-x_{k+1}\|^{2}=L^{2}\alpha^{2}\|v_k\|^{2}.
\tag{13}
\]

Taking expectation over $i_k$ and using \eqref{11},
\[
\E\!\big[\Delta_{k+1}\mid\mathcal{F}_k\big]
\le\left(1-\frac{1}{n}\right)\Delta_k
+\frac{1}{n}L^{2}\alpha^{2}\E\!\big[\|v_k\|^{2}\mid\mathcal{F}_k\big].
\tag{14}
\]

Insert bound \eqref{11} into \eqref{14}:
\[
\E\!\big[\Delta_{k+1}\mid\mathcal{F}_k\big]
\le\left(1-\frac{1}{n}\right)\Delta_k
+\frac{2L^{2}\alpha^{2}}{n}\big(\|\nabla f(x_k)\|^{2}+\Delta_k\big).
\tag{15}
\]

Collect the $\Delta_k$ terms:
\[
\E\!\big[\Delta_{k+1}\mid\mathcal{F}_k\big]
\le\Big(1-\frac{1}{n}+\frac{2L^{2}\alpha^{2}}{n}\Big)\Delta_k
+\frac{2L^{2}\alpha^{2}}{n}\|\nabla f(x_k)\|^{2}.
\tag{16}
\]

\item[Step 9.] \textbf{Define the Lyapunov function.}
Let
\[
\Phi_k:= f(x_k)+c\,\Delta_k,\qquad
c:=\frac{\alpha L}{1-\alpha L}>0,
\tag{17}
\]
where the choice of $c$ will cancel the $\|\nabla f(x_k)\|^{2}$ terms.

\item[Step 10.] \textbf{One‑step Lyapunov decrease.}
Take total expectation of \eqref{12} and add $c$ times \eqref{16}:
\begin{align*}
\E[\Phi_{k+1}]
&\le \E[f(x_k)]-\alpha(1-L\alpha)\E[\|\nabla f(x_k)\|^{2}]
+L\alpha^{2}\E[\Delta_k] \\
&\quad +c\Bigg\{\Big(1-\frac{1}{n}+\frac{2L^{2}\alpha^{2}}{n}\Big)\E[\Delta_k]
+\frac{2L^{2}\alpha^{2}}{n}\E[\|\nabla f(x_k)\|^{2}]\Bigg\}.
\end{align*}
Group the $\|\nabla f(x_k)\|^{2}$ coefficients:
\[
\bigg[-\alpha(1-L\alpha)+\frac{2cL^{2}\alpha^{2}}{n}\bigg]\E[\|\nabla f(x_k)\|^{2}]
\]
and the $\Delta_k$ coefficients:
\[
\bigg[L\alpha^{2}+c\Big(1-\frac{1}{n}+\frac{2L^{2}\alpha^{2}}{n}\Big)\bigg]\E[\Delta_k].
\]

\item[Step 11.] \textbf{Choose $c$ to eliminate the gradient term.}
Set
\[
-\alpha(1-L\alpha)+\frac{2cL^{2}\alpha^{2}}{n}= -\frac{\alpha}{2},
\]
which is satisfied when $c=\frac{\alpha L}{1-\alpha L}$ (recall $n\ge1$ and $\alpha\le\frac{1}{3L}$). With this $c$ we obtain
\[
\E[\Phi_{k+1}]
\le \E[\Phi_k]-\frac{\alpha}{2}\E[\|\nabla f(x_k)\|^{2}]
+\underbrace{\Big(L\alpha^{2}+c\Big(1-\frac{1}{n}+\frac{2L^{2}\alpha^{2}}{n}\Big)\Big)}_{\le 0}\E[\Delta_k].
\]
The term in braces is non‑positive for $\alpha\le1/(3L)$ (simple algebraic verification), hence it can be dropped:
\[
\E[\Phi_{k+1}]
\le \E[\Phi_k]-\frac{\alpha}{2}\E[\|\nabla f(x_k)\|^{2}].
\tag{18}
\]

\item[Step 12.] \textbf{Summation over $k$.}
Sum \eqref{18} for $k=0,\dots,T-1$ and telescope:
\[
\Phi_0-\E[\Phi_T]\ge \frac{\alpha}{2}\sum_{k=0}^{T-1}\E[\|\nabla f(x_k)\|^{2}].
\]
Since $\Phi_T\ge f^{\star}$,
\[
\sum_{k=0}^{T-1}\E[\|\nabla f(x_k)\|^{2}]
\le \frac{2}{\alpha}\big(\Phi_0-f^{\star}\big)
= \frac{2}{\alpha}\big(f(x_0)-f^{\star}+c\Delta_0\big).
\tag{19}
\]

\item[Step 13.] \textbf{Divide by $T$ and substitute $c$.}
Using $c=\frac{\alpha L}{1-\alpha L}$,
\[
\frac{1}{T}\sum_{k=0}^{T-1}\E[\|\nabla f(x_k)\|^{2}]
\le \frac{2\big(f(x_0)-f^{\star}\big)}{\alpha T}
+\frac{2c\Delta_0}{\alpha T}
= \frac{2\big(f(x_0)-f^{\star}\big)}{\alpha T}
+\frac{2L}{1-\alpha L}\cdot\frac{\Delta_0}{T}.
\]
Multiplying numerator and denominator of the second term by $\alpha$ yields the bound stated in \eqref{4}:
\[
\frac{1}{T}\sum_{k=0}^{T-1}\E[\|\nabla f(x_k)\|^{2}]
\le
\frac{2\big(f(x_0)-f^{\star}\big)}{\alpha T}
+\frac{4L\alpha}{1-\alpha L}\,\E[\Delta_0].
\qquad\blacksquare
\]

\end{enumerate}
%%%%%%%%%%%%%%%%%%%%%%%%%%%%%%%%%%%%%%%%%%%%%%%%%%%%%%%%%%%%

\section*{Rate Derivation (With Constants)}
\addcontentsline{toc}{section}{Rate Derivation (With Constants)}
The bound \eqref{4} contains two terms:
\[
\underbrace{\frac{2\big(f(x_0)-f^{\star}\big)}{\alpha T}}_{\text{optimization term}}
\qquad\text{and}\qquad
\underbrace{\frac{4L\alpha}{1-\alpha L}\,\E[\Delta_0]}_{\text{gradient‑table term}}.
\]
Choosing $\alpha = \min\big\{\frac{1}{3L},\,\sqrt{\frac{2\big(f(x_0)-f^{\star}\big)}{4L\,\E[\Delta_0]\,T}}\big\}$
balances the two contributions and yields
\[
\frac{1}{T}\sum_{k=0}^{T-1}\E[\|\nabla f(x_k)\|^{2}]
= \mathcal{O}\!\Big(\frac{1}{\sqrt{T}}\Big).
\]
If the stored gradients are initially accurate (i.e., $\E[\Delta_0]=0$), the second term vanishes and we obtain the faster $\mathcal{O}\!\big(\frac{1}{\alpha T}\big)$ rate, i.e. linear decay in $1/T$.

%%%%%%%%%%%%%%%%%%%%%%%%%%%%%%%%%%%%%%%%%%%%%%%%%%%%%%%%%%%%

\section*{Special Cases}
\addcontentsline{toc}{section}{Special Cases}
\begin{enumerate}
\item \textbf{Convex $f$.} When $f$ is convex, the same proof yields the classical $\mathcal{O}(1/T)$ convergence of the function values $f(\bar x_T)-f^{\star}$, because $\|\nabla f(x_k)\|^{2}$ lower‑bounds the suboptimality via strong smoothness.

\item \textbf{Single component ($n=1$).} SAGA reduces to standard gradient descent with stepsize $\alpha\le 1/(3L)$; the bound collapses to the familiar $\frac{2(f(x_0)-f^{\star})}{\alpha T}$.

\item \textbf{Large $n$ limit.} As $n\to\infty$, the table error $\Delta_k$ decays faster (the factor $1-\frac{1}{n}$ in \eqref{16} approaches $1$), and the algorithm behaves similarly to SGD but with a variance reduction term of order $\alpha^{2}$.
\end{enumerate}

%%%%%%%%%%%%%%%%%%%%%%%%%%%%%%%%%%%%%%%%%%%%%%%%%%%%%%%%%%%%

\end{document}